\Needspace{7\baselineskip}
\begin{proof}[Proof of Proposition~\ref{prop:rewrite-uncapped-complements-bounds}]
\phantomsection
\label{proof:rewrite-uncapped-complements-bounds}
Equations~\eqref{eq:rewrite-uncapped-complements-z-bound}
and~\eqref{eq:rewrite-uncapped-complements-y-bound} follow because
$\sigma/(\sigma-1)<0$. Define $k=K/(AL)$. The resource constraint
\eqref{eq:rewrite-eq-kdot}, nonnegative consumption and compute expenditure,
and $L=N$ imply
\begin{equation}
 \dot k\leq
 \omega_L^{\sigma(1-\alpha)/(\sigma-1)}k^\alpha
 -(\delta+n+\gamma)k.
 \label{eq:rewrite-uncapped-complements-k-bound}
\end{equation}
Since $\delta+n+\gamma>0$ and $0<\alpha<1$, the right-hand side is
negative above a finite positive level. In particular,
\[
 k_t\leq
 \max\left\{k_0,
 \left[\frac{\omega_L^{\sigma(1-\alpha)/(\sigma-1)}}
 {\delta+n+\gamma}\right]^{1/(1-\alpha)}\right\}.
\]
Equation~\eqref{eq:rewrite-uncapped-complements-y-bound} then bounds
$Y/(AL)$ uniformly. Since $A_t=A_0e^{\gamma t}$ and $L=N$,
$Y/N$ is bounded above by a constant times $e^{\gamma t}$.
Taking logarithms, dividing by $t$, and taking the upper limit gives
\eqref{eq:rewrite-uncapped-complements-growth-bound}.
\end{proof}

% Additional results retained for possible later use; not part of the proposition.
% The same bounds keep $K$ and $Y$ finite on every bounded time interval.
% For any such horizon $T$, integrating the resource constraint gives
% \[
%  \int_0^T M_t\dd t\leq K_0+\int_0^T Y_t\dd t<\infty.
% \]
% The uncapped RSI technology and H\"older's inequality imply
% \begin{align}
%  B_T^{1-\eta}
%  &=B_0^{1-\eta}+(1-\eta)\chi\int_0^T M_t^\eta\dd t
%  \nonumber\\
%  &\leq B_0^{1-\eta}
%  +(1-\eta)\chi T^{1-\eta}
%  \left(\int_0^T M_t\dd t\right)^\eta<\infty.
%  \label{eq:rewrite-uncapped-complements-b-bound}
% \end{align}
% Thus $B$ cannot diverge at a finite date either. This argument bounds
% stocks and output; it is not a regularity theorem for all instantaneous
% control choices or an equilibrium-existence proof.
%
% There is also no finite-horizon unbounded-payoff argument of the kind used
% for substitution in Proposition~\ref{prop:rewrite-research-scale}.
% At given positive continuous paths of $K,A,L$ on $[0,T]$, optimized operating
% profit satisfies
% \[
%  0\leq\pi_t(B)\leq p_XX\leq(1-\alpha)Y
%  \leq(1-\alpha)\omega_L^{\sigma(1-\alpha)/(\sigma-1)}
%  K_t^\alpha(A_tL_t)^{1-\alpha}.
% \]
% This upper bound is independent of the developer's choice of $B$.
% For a given locally integrable $r$, the discount factor
% $D_t=\exp(-\int_0^t r_s\dd s)$ has a positive minimum on $[0,T]$.
% Discounted operating revenue is bounded, whereas the cost of research
% tends to infinity with $\int_0^T M_t\dd t$. The developer's truncated
% supremum is therefore finite, and net payoff tends to minus infinity as
% total research expenditure diverges. This conclusion holds for every
% $0<\eta<1$, but does not prove that an infinite-horizon optimum exists.

\Needspace{7\baselineskip}
\begin{proof}[Proof of Proposition~\ref{cond:rewrite-uncapped-complements-limits}]
\phantomsection
\label{proof:rewrite-uncapped-complements-limits}
All auxiliary variables below are confined to this proof. Write
$x=X/(AL)$, $k=K/(AL)$, $c=C/(AL)$, $y=Y/(AL)$, $u=U/(AL)$,
$m=M/(AL)$, $d=\rho-n>0$, and $\nu=\delta+n+\gamma$.
No convergence of a growth rate or allocation ratio is assumed.

\textit{Step 1. Static optimality gives a limiting production function.}
Put $\varphi=(\sigma-1)/\sigma<0$ and
$f(x)=[\omega_L+\omega_Xx^\varphi]^{(1-\alpha)/\varphi}$, so $y=k^\alpha f(x)$.
The monopoly condition~\eqref{eq:rewrite-monopoly-foc} requires
$e_X=(1-s_X)/\sigma+\alpha s_X<1$. Thus
\begin{align}
 s_X>s_\infty&:=\frac{1-\sigma}{1-\alpha\sigma},&
 0<x<x_\infty&:=
 \left[\frac{\omega_Ls_\infty}{\omega_X(1-s_\infty)}\right]^{1/\varphi}.
 \label{eq:rewrite-uncapped-complements-x-limit}
\end{align}
Define $f_\infty=f(x_\infty)>0$.
In particular, $y\leq f_\infty k^\alpha$ and $0<u\leq x_\infty/B$.

The monopoly condition can be written
\[
 k^\alpha v(x)=1/B,\qquad
 v(x)=\frac{(1-\alpha)s_Xf(x)(1-e_X)}x.
\]
On $(0,x_\infty)$, $v$ is strictly decreasing, from infinity to zero.
Indeed, with a prime here denoting differentiation with respect to $\log x$,
$e_X'=(\alpha-\sigma^{-1})\varphi s_X(1-s_X)>0$ and
$-\mathrm d\log v/\mathrm d\log x=e_X+e_X'/(1-e_X)>\alpha s_X$.
Consequently $x$ increases with $B$ and $k$, while
\[
 \frac{\partial\log y}{\partial\log k}
 =\alpha+\frac{\alpha(1-\alpha)s_X}{e_X+e_X'/(1-e_X)}<1.
\]
Thus $y/k$ decreases with $k$ and increases with $B$.
For $k$ in any compact interval bounded away from zero, as $B\to\infty$,
\[
 x\to x_\infty,\qquad y\to f_\infty k^\alpha,\qquad u\to0
\]
uniformly. These are properties of the monopoly allocation at given $k,B$,
not assumptions about the equilibrium path. They also imply that, for any
prescribed finite return, sufficiently small $k$ and sufficiently large $B$
make $r$ exceed that return, uniformly over those small values of $k$.

\textit{Step 2. Household and developer optimality bound consumption.}
Proposition~\ref{prop:rewrite-uncapped-complements-bounds} bounds $k$ and
$y$ above by finite constants $k_{\max},y_{\max}$.
The household Euler equation and $r\geq-\delta$ give
$\dot c/c\geq-(\delta+\rho+\gamma)$.
Integrating the resource constraint over $[t,t+1]$ therefore gives
\[
 c_t\int_0^1 e^{-(\delta+\rho+\gamma)v}\dd v
 \leq\int_t^{t+1}c_s\dd s\leq k_{\max}+y_{\max}.
\]
Hence $c$ has a finite uniform upper bound $c_{\max}$.

Write $D_{t,s}=\exp(-\int_t^s r_v\dd v)$.
Euler implies $\int_t^\infty D_{t,s}C_s\dd s=C_t/d$.
Integrating the household budget~\eqref{eq:rewrite-household-budget}
and using its TVC~\eqref{eq:rewrite-household-tvc} shows that
the discounted integral of labor income plus developer net profit converges.
Since labor income is nonnegative, the discounted integral of developer net
profit has a limit that is either finite or minus infinity. The latter is impossible:
the developer can stop research and obtain nonnegative operating profits.
By continuation optimality, the developer's continuation value
$V_t=\int_t^\infty D_{t,s}\Pi_s\dd s$ is therefore finite and nonnegative.
It follows that
\begin{equation}
 \frac{C_t}{d}
 =K_t+V_t+\int_t^\infty D_{t,s}w_sL_s\dd s,\qquad
 c_t\geq dk_t,\qquad
 \int_t^\infty D_{t,s}w_sL_s\dd s\leq\frac{C_t}{d}.
 \label{eq:rewrite-complements-household-wealth-bound}
\end{equation}
Here $V_t$ is an auxiliary continuation value, not an additional traded
asset or a change to the household budget.

\textit{Step 3. Discounted research spending vanishes relative to effective labor.}
Fix a positive rate $a$. Step 1 permits a small $k_a>0$ and a sufficiently
late date such that $r-(n+\gamma)\geq a$ whenever $k\leq k_a$.
On $[k_a,k_{\max}]$, the wage equation~\eqref{eq:rewrite-wage} and the
static monopoly allocation give a uniform positive lower bound $\ell_a$ on $w/A$.
Since $r\geq-\delta$, there is a finite constant $H_a>0$ such that
\[
 a\leq r-(n+\gamma)+H_a w/A
\]
at every sufficiently late date. Multiplication by $D_{t,s}A_sL_s$,
integration, and~\eqref{eq:rewrite-complements-household-wealth-bound} give
\begin{equation}
 \int_t^\infty D_{t,s}A_sL_s\dd s
 \leq\frac{A_tL_t+H_aC_t/d}{a}
 \leq J A_tL_t
 \label{eq:rewrite-complements-effective-labor-pv-bound}
\end{equation}
for a fixed finite $J$. The omitted terminal contribution from discounted
effective labor is nonpositive; no discount-rate limit is used.

The research condition~\eqref{eq:rewrite-research-foc} and costate
equation~\eqref{eq:rewrite-costate} imply
\begin{equation}
 \frac{\mathrm d(qB)}{\mathrm dt}
 =rqB-U+\frac{1-\eta}{\eta}M.
 \label{eq:rewrite-complements-shadow-stock}
\end{equation}
Since $B$ is nondecreasing, $U_s\leq x_\infty A_sL_s/B_t$ for $s\geq t$.
Integrating~\eqref{eq:rewrite-complements-shadow-stock} first over finite
intervals and then using the developer TVC~\eqref{eq:rewrite-developer-tvc} gives
\begin{equation}
 q_tB_t+\frac{1-\eta}{\eta}
       \int_t^\infty D_{t,s}M_s\dd s
 =\int_t^\infty D_{t,s}U_s\dd s
 \leq\frac{x_\infty J}{B_t}A_tL_t.
 \label{eq:rewrite-complements-research-pv-bound}
\end{equation}
The right side is finite by~\eqref{eq:rewrite-complements-effective-labor-pv-bound};
nonnegativity then ensures that the research integral is finite as well.
Its value divided by $A_tL_t$ tends to zero. This argument requires
$0<\eta<1$, but neither $\eta\leq1/2$ nor $\eta<\alpha$.

\textit{Step 4. Capital cannot repeatedly approach zero.}
Let $F(k)=f_\infty k^\alpha-\nu k$.
Choose $\varepsilon>0$ small enough that $F$ is positive and increasing on
$(0,\varepsilon]$. Step 1 also permits a sufficiently late date and
$a>0$ such that $\dot c/c=r-\rho-\gamma\geq a$ whenever $k\leq\varepsilon$.
For any $j\leq\varepsilon$, the equilibrium path cannot enter
$\{k\leq j,\ c\geq F(j)\}$: there $\dot k\leq F(k)-c\leq0$,
whereas $c$ continues to grow at least at rate $a$, contradicting its upper bound.

If $k$ stayed below $\varepsilon$ forever, consumption would likewise
exceed its upper bound. Hence the path eventually reaches $\varepsilon$.
Suppose arbitrarily late crossings from $\varepsilon$ to $\varepsilon/2$
were possible. Let $t_0$ be the last crossing of $\varepsilon$ before the
first subsequent crossing $t_1$ of $\varepsilon/2$.
Throughout this interval $k\in[\varepsilon/2,\varepsilon]$ and $c$ increases.
Equation~\eqref{eq:rewrite-complements-household-wealth-bound} gives
$c_{t_0}\geq d\varepsilon$, so
\[
 t_1-t_0\leq T_\varepsilon
 :=a^{-1}\log\!\left(\frac{c_{\max}}{d\varepsilon}\right)<\infty.
\]
At $t_1$, the preceding invariant-region argument requires
$c_{t_1}<F(\varepsilon/2)$. Thus $c_s<F(\varepsilon/2)$ throughout the crossing.
The uniform convergence in Step 1, applied on the crossing interval, yields
numbers $\epsilon_{t_0}\to0$ such that
\[
 \dot k_s
 =y_s-u_s-\nu k_s-c_s-m_s
 \geq-\epsilon_{t_0}-m_s.
\]
On this interval $r$ is bounded above by a fixed positive $R_\varepsilon$.
Since $n+\gamma>0$, equation~\eqref{eq:rewrite-complements-research-pv-bound}
then implies
\[
 \frac{\varepsilon}{2}
 \leq T_\varepsilon\epsilon_{t_0}+\int_{t_0}^{t_1}m_s\dd s
 \leq T_\varepsilon\epsilon_{t_0}
 +e^{R_\varepsilon T_\varepsilon}
   \frac{\eta x_\infty J}{(1-\eta)B_{t_0}}
 \longrightarrow0,
\]
a contradiction. Therefore $k$ is eventually bounded away from zero,
and so is $c\geq dk$. In particular, $r$ is now bounded above.

\textit{Step 5. Research per unit of effective labor vanishes and Ramsey dynamics determine the limit.}
Differentiating the research condition and using the costate equation gives
\begin{equation}
 (1-\eta)g_M
 =r-\frac{\eta\chi UB^{\eta-1}}{M^{1-\eta}}\leq r.
 \label{eq:rewrite-complements-research-growth-bound}
\end{equation}
Choose a positive $R$ above all eventual values of $r$ and put $H=R/(1-\eta)$.
Then $M_s\geq M_t e^{-H(t-s)}$ on $[t-1,t]$ at sufficiently late dates.
Consequently
\[
 \frac1{A_{t-1}L_{t-1}}
 \int_{t-1}^t D_{t-1,s}M_s\dd s
 \geq e^{n+\gamma-R-H}m_t.
\]
The left side tends to zero by~\eqref{eq:rewrite-complements-research-pv-bound}.
Hence $m\to0$, and Step 1 already gives $u\to0$.

The normalized equilibrium equations therefore approach, uniformly on
the compact positive region containing the path, the Ramsey system
\begin{equation}
 \dot k=F(k)-c,\qquad
 \dot c=c\,[\alpha f_\infty k^{\alpha-1}-\delta-\rho-\gamma].
 \label{eq:rewrite-complements-limiting-ramsey}
\end{equation}
For clarity, the following argument verifies convergence rather than
inferring it solely from convergence of the equations.
This limiting system has one positive stationary point,
\[
 k_\infty=
 \left[\frac{\alpha f_\infty}{\rho+\gamma+\delta}\right]^{1/(1-\alpha)},
 \qquad c_\infty=F(k_\infty)>0.
\]
Its Jacobian has negative determinant, so it is a saddle.
The standard Ramsey phase portrait is discussed by \citet{groth2013ramsey}.
There is no nonconstant complete orbit confined to a compact subset of
the positive quadrant. To see this, the multiplier $1/c^2$ gives
\[
 \frac{\partial}{\partial k}\frac{F(k)-c}{c^2}
 +\frac{\partial}{\partial c}
   \frac{\alpha f_\infty k^{\alpha-1}-\delta-\rho-\gamma}{c}
 =\frac{\rho-n}{c^2}>0.
\]
The Bendixson--Dulac argument excludes both periodic orbits and homoclinic
loops. The planar limit-set theorem, the unique stationary point, and its
saddle character then exclude a nonconstant complete orbit confined to
such a compact set: its forward and backward limits would otherwise
require a periodic orbit or a connection back to that same saddle.

For any sequence of dates tending to infinity, time translations of our
bounded positive path have a subsequence converging on every bounded
time interval to a complete orbit of
\eqref{eq:rewrite-complements-limiting-ramsey}. This follows from bounded
derivatives, uniform convergence of the vector fields, and the integral
form of the differential equations. The preceding paragraph forces that
orbit to be the stationary point. Thus every accumulation point is
$(k_\infty,c_\infty)$, proving convergence of the original path.

It follows that $y\to y_\infty=f_\infty k_\infty^\alpha>0$ and
\begin{equation}
 r\to\rho+\gamma,\qquad
 \frac KY\to\frac{\alpha}{\rho+\gamma+\delta},\qquad
 \frac CY\to1-\frac{\alpha(\delta+n+\gamma)}{\rho+\gamma+\delta}>0.
 \label{eq:rewrite-complements-euler-limit}
\end{equation}
Positivity follows because the numerator of the final expression is
$\rho-n+(1-\alpha)(\delta+n+\gamma)>0$.
This proves~\eqref{eq:rewrite-uncapped-complements-allocation-ratios}.
Since $u,m\to0$ and $y\to y_\infty>0$, it follows that
$U/Y=u/y\to0$ and $M/Y=m/y\to0$, proving
\eqref{eq:rewrite-uncapped-complements-compute-shares}.
It also gives $\dot k\to0$ and $g_C\to n+\gamma$.

\textit{Step 6. Output and wage growth converge as well.}
A convergent level need not have a convergent derivative.
The bound on $g_M$ from Step 5 and the RSI technology give
\[
 B_t^{1-\eta}
 \geq(1-\eta)\chi M_t^\eta\int_0^1e^{-\eta H v}\dd v.
\]
Thus $g_B=\chi M^\eta/B^{1-\eta}$ is bounded above.
At $x_\infty$, the marginal-revenue function $v$ defined in Step 1 has
a simple zero: $v(x_\infty)=0$ and $v'(x_\infty)<0$, where the prime
now denotes differentiation with respect to $x$.
Differentiating $k^\alpha v(x)=1/B$ yields
\[
 \dot x=-\frac{v(x)}{v'(x)}
 \left(g_B+\alpha\frac{\dot k}{k}\right)\longrightarrow0.
\]
Hence $\dot y\to0$ and $\dot s_X\to0$.
Since $y_\infty>0$ and $1-s_\infty>0$, $g_{Y/N}\to\gamma$ and,
by~\eqref{eq:rewrite-wage}, $g_w\to\gamma$.
The limiting income shares follow by substituting $s_\infty$ into
\eqref{eq:rewrite-wage} and~\eqref{eq:rewrite-ai-price}.
The result is conditional on equilibrium existence and $B\to\infty$,
not an equilibrium-existence theorem or an assumption of balanced growth.
\end{proof}
